\documentclass{beamer}
\usepackage[utf8]{inputenc}
\usepackage{booktabs}
\usepackage{graphicx}
\usepackage{hyperref}
\usetheme{default}
\usecolortheme{default}
\setbeamertemplate{navigation symbols}{}
\setbeamercolor{background canvas}{bg=white}

\title{Theory and Model Comparison}
\author{Deepfake‑ViT Project}
\date{\today}

\begin{document}

\frame{\titlepage}

% -------------------------------------------------
\section{Theoretical Foundations}
\begin{frame}{Deepfake Forensics Goal}
  \begin{itemize}
    \item Detect subtle spatial, spectral, and semantic artifacts.
    \item Benchmark Vision Transformers (ViT) vs. Convolutional Neural Networks (CNN).
  \end{itemize}
\end{frame}

% -------------------------------------------------
\section{Backbone Architectures}
\subsection{Meta DINOv3 ViT‑S/16}
\begin{frame}{ViT‑S/16 Overview}
  \begin{itemize}
    \item 28.69M parameters, patch size 16×16, 12 encoder layers, dim 384, SwiGLU Gated MLP (strictly matched to ConvNeXt-Tiny 28.12M).
    \item Global all‑to‑all self‑attention → captures whole‑face illumination and symmetry.
    \item Strengths: Global receptive field, high scalability with pretrained foundation.
    \item Weaknesses: No spatial inductive bias, quadratic cost $O(N^2 D)$.
  \end{itemize}
\end{frame}

\subsection{Meta DINOv3 ConvNeXt‑Tiny}
\begin{frame}{ConvNeXt‑Tiny Overview}
  \begin{itemize}
    \item 28.1M parameters, 4 hierarchical stages, channels [96,192,384,768].
    \item 7×7 depthwise kernels + LayerScale → retains 2D locality.
    \item Strengths: Fast inference (153 FPS), excellent on local boundary artifacts.
    \item Weaknesses: Limited long‑range context, receptive field bound by depth.
  \end{itemize}
\end{frame}

\subsection{Classical CNN Baselines}
\begin{frame}{ResNet‑50 & EfficientNet‑B4}
  \begin{itemize}
    \item Proven convergence, small memory footprint.
    \item Limitations: Small $3\times3$ kernels, BatchNorm sensitivity to domain shift.
  \end{itemize}
\end{frame}

% -------------------------------------------------
\section{Joint Weighted Ensemble}
\begin{frame}{Ensemble Formula}
  \[ P_{\text{ens}} = 0.65\,P_{\text{ViT}} + 0.35\,P_{\text{CNN}} \]
  \begin{itemize}
    \item Improves accuracy to 97.88 % and ROC‑AUC to 99.74 %.
    \item Reduces false negatives by 12.9 % over ViT alone.
  \end{itemize}
\end{frame}

% -------------------------------------------------
\section{Model Comparison Matrix}
\begin{frame}{Key Evaluation Aspects}
  \begin{tabular}{lcccc}
    \toprule
    Aspect & ResNet‑50 & ConvNeXt‑Tiny & ViT‑S/16 & Ensemble \\
    \midrule
    Inductive Bias & Strong & Strong & None & Dual \\
    Information Scope & Local & Mid‑Local & Global & Dual \\
    Data Requirement & Moderate & Moderate & High & High \\
    Computational Cost & $O(C K^2 HW)$ & $O(C K^2 HW)$ & $O(N^2 D)$ & Sum \\
    Best For & Small datasets & Real‑time & Large pre‑train & Mission‑critical \\
    \bottomrule
  \end{tabular}
\end{frame}

% -------------------------------------------------
\section{When ViT > CNN}
\begin{frame}{ViT‑Preferred Scenarios}
  \begin{enumerate}
    \item Diffusion‑based generators (global illumination inconsistencies).
    \item Whole‑image GAN synthesis (subtle geometric asymmetries).
    \item Availability of large pretrained foundations.
  \end{enumerate}
\end{frame}

% -------------------------------------------------
\section{When CNN > ViT}
\begin{frame}{CNN‑Preferred Scenarios}
  \begin{enumerate}
    \item FaceSwap / local blending seams (boundary artifacts).
    \item Resource‑constrained or real‑time video pipelines.
    \item Training from scratch on limited data (<10k samples).
  \end{enumerate}
\end{frame}

% -------------------------------------------------
\section{Practical Guidance}
\begin{frame}{Decision Guide}
  \begin{itemize}
    \item <5k samples → ResNet / ConvNeXt.
    \item 20k–50k samples → ConvNeXt‑Tiny.
    \item >100k samples → ViT‑S/16.
    \item Accuracy‑critical → Joint Ensemble.
    \item Edge devices → ConvNeXt‑Tiny.
  \end{itemize}
\end{frame}

\end{document}
